\section{Convergence and diagnostics}
\label{sec:convergence}

A multireference calculation is only as good as the stationary point
it was evaluated at, and the failure modes here are quiet ones.  This
chapter collects what the package reports and how to read it.

\subsection{The criterion is the gradient, not the energy}

The orbital optimization returns both a point and the gradient norm at
that point, and the norm is the verdict.  The energy alone cannot
serve: this project measured a run in which the energy kept descending
while the parameters walked to magnitudes of $10^{13}$ --- the model
had no stationary point in the direction it was sliding along, and an
energy-only stopping rule called it a success.  Treat a convergence
claim without a gradient norm as no claim at all.

The stopping reason distinguishes the ways a run can end:

\begin{itemize}
  \item \texttt{CONVERGED}: the gradient reached the refinement
    target.  This is the only reason that is a stationarity claim.
  \item \texttt{STALLED\_GRADIENT}: the iteration stopped because a
    step no longer reduced the gradient.  The point may still be
    useful, but it is a candidate, not a certificate.
  \item \texttt{GUARD\_STOPPED}: no damped step lowered the energy.
    On a converged point this is the normal way to finish a Newton
    pass.
  \item \texttt{ITERATION\_CAP}: the budget ran out.  Re-check the
    gradient before using the point.
\end{itemize}

\subsection{Multiple stationary points}

The CASSCF energy is not convex in the orbital parameters.  Different
starting points can reach different stationary points whose energies
differ by more than a bond energy; the water molecule at 6-31G has two
minima separated by 13~kcal/mol from the same molecule and basis.
Nothing in the eigensolver signals this: both points are stationary.

Practical consequences:

\begin{itemize}
  \item Report the stationary-point spread alongside any
    multireference number, or run several starts and keep the lowest.
  \item A point with zero gradient can still be a saddle.  The kernel
    exposes a curvature check over the Hessian
    (\texttt{stationary.curvature\_report}); when a number matters,
    run it.  The orbitals obtained straight from the reference SCF
    ($\theta = 0$) are, for the water example in this tutorial, a
    saddle rather than a minimum.
  \item Treat the spread as protocol, not error: fixed settings can
    still select different stationary points on different machines,
    and the package reports the choice it made rather than
    normalizing it away.
\end{itemize}

\subsection{The perturbation layer is the more sensitive one}

NEVPT2 depends on the orbitals more sensitively than the CASSCF energy
does.  Convergence gates set for the energy can leave the correction
unstable: this project measured orbital-gradient stalls that moved the
correction by several kcal/mol between runs whose energies agreed to
the microhartree.  When the target quantity is the correction,

\begin{itemize}
  \item tighten the orbital-gradient gate rather than the energy
    gate, and read the gradient norm the run reports;
  \item confirm the correction's stability across independent runs
    (different starting parameters, or the two state-averaging passes
    of Chapter~\ref{sec:state-averaged}), instead of trusting a
    single converged number.
\end{itemize}
